\documentclass[12pt]{article}
\usepackage{amsmath, amssymb}
\usepackage[margin=1.7cm]{geometry}
\usepackage{hyperref}
\usepackage{graphicx}
\usepackage{tikz}
\usepackage{enumitem}
\setlength{\parindent}{0pt}
\setlist{itemsep=0.35em, topsep=0.35em}

\title{Äquivalenzumformungen mit Brüchen}
\author{}
\date{14.02.2026}

\begin{document}
\maketitle

\section*{Wiederholung: Brüche vereinfachen}

\subsection*{Brüche kürzen}
Kürze die Brüche soweit wie möglich.
\begin{enumerate}
  \item \(\frac{4}{8}\)
  \item \(\frac{6}{9}\)
  \item \(\frac{12}{18}\)
  \item \(\frac{15}{25}\)
  \item \(\frac{24}{36}\)
\end{enumerate}

\subsection*{Gleichnamige Brüche addieren/subtrahieren}
Rechne aus. Schreibe jeden Schritt.
\begin{enumerate}
  \item \(\frac{1}{4} + \frac{2}{4}\)
  \item \(\frac{3}{5} - \frac{1}{5}\)
  \item \(\frac{2}{7} + \frac{4}{7}\)
  \item \(\frac{5}{6} - \frac{3}{6}\)
\end{enumerate}

\subsection*{Ungleichnamige Brüche}
Rechne aus. Finde zuerst den Hauptnenner.
\begin{enumerate}
  \item \(\frac{1}{2} + \frac{1}{4}\)
  \item \(\frac{2}{3} + \frac{1}{6}\)
  \item \(\frac{1}{3} - \frac{1}{6}\)
  \item \(\frac{3}{4} + \frac{1}{2}\)
  \item \(\frac{2}{5} + \frac{3}{10}\)
\end{enumerate}

\section*{Äquivalenzumformungen mit Brüchen}

\subsection*{Wichtige Grundregel}
Multiplikation mit dem Kehrwert:
\[
  \frac{a}{b} = c \quad \xrightarrow{\cdot \frac{b}{a}} \quad 1 = c \cdot \frac{b}{a}
\]

\subsection*{Aufgabe 1: Einfache Gleichungen mit Brüchen}
Löse nach der Variable.
\begin{enumerate}
  \item \(\frac{x}{2} = 5\)
  \item \(\frac{y}{3} = 4\)
  \item \(\frac{a}{5} = 2\)
  \item \(\frac{z}{4} = 7\)
  \item \(-\frac{x}{2} = 3\)
  \item \(\frac{y}{-3} = 5\)
\end{enumerate}

\subsection*{Aufgabe 2: Bruch umformen}
Bring die Gleichung in die Form \(x = \text{Zahl}\).
\begin{enumerate}
  \item \(\frac{2x}{3} = 4\)
  \item \(\frac{3y}{4} = 6\)
  \item \(\frac{5a}{2} = 10\)
  \item \(\frac{4z}{5} = 8\)
\end{enumerate}

\subsection*{Aufgabe 3: Bruch isolieren}
Löse nach \(x\).
\begin{enumerate}
  \item \(\frac{x + 1}{2} = 3\)
  \item \(\frac{2x + 3}{4} = 5\)
  \item \(\frac{3x - 1}{5} = 4\)
  \item \(\frac{4x + 2}{6} = 3\)
\end{enumerate}

\section*{Gleichungen mit mehreren Brüchen}

\subsection*{Neuer Trick: Hauptnenner finden}
Multipliziere beide Seiten mit dem Hauptnenner (kleinstes gemeinsames Vielfaches aller Nenner).

\subsection*{Aufgabe 4: Zwei Brüche auf einer Seite}
Löse die Gleichungen.
\begin{enumerate}
  \item \(\frac{x}{2} + \frac{x}{4} = 6\)
  \item \(\frac{y}{3} + \frac{2y}{6} = 8\)
  \item \(\frac{a}{5} + \frac{2a}{5} = 9\)
  \item \(\frac{z}{4} + \frac{3z}{8} = 5\)
\end{enumerate}

\subsection*{Aufgabe 5: Brüche auf beiden Seiten}
Löse die Gleichungen.
\begin{enumerate}
  \item \(\frac{x}{2} + \frac{x}{3} = 5\)
  \item \(\frac{y}{4} - \frac{y}{6} = 3\)
  \item \(\frac{a}{3} + \frac{a}{5} = 8\)
  \item \(\frac{2z}{5} - \frac{z}{10} = 6\)
\end{enumerate}

\subsection*{Aufgabe 6: Brüche mit Konstanten}
Löse die Gleichungen.
\begin{enumerate}
  \item \(\frac{x}{2} + 3 = 7\)
  \item \(\frac{y}{3} - 2 = 4\)
  \item \(\frac{a}{4} + 5 = 9\)
  \item \(\frac{z}{5} - 1 = 6\)
\end{enumerate}

\section*{Brüche beseitigen}

\subsection*{Schritt 1: Hauptnenner finden}
\subsection*{Schritt 2: Multiplizieren}
\subsection*{Schritt 3: Auflösen}

\subsection*{Aufgabe 7: Schrittweise vorgehen}
Löse die Gleichungen. Zeige jeden Schritt:
\begin{enumerate}
  \item \(\frac{x + 1}{2} + \frac{x}{4} = 5\)
  \item \(\frac{2y}{3} - \frac{y - 1}{6} = 4\)
  \item \(\frac{3a + 1}{4} + \frac{a}{8} = 6\)
  \item \(\frac{4z - 2}{5} - \frac{z}{10} = 3\)
\end{enumerate}

\subsection*{Aufgabe 8: Komplexere Gleichungen}
\begin{enumerate}
  \item \(\frac{2x + 1}{3} + \frac{x - 2}{6} = 4\)
  \item \(\frac{3y - 2}{4} - \frac{y + 1}{8} = 2\)
  \item \(\frac{5a + 1}{5} - \frac{2a - 3}{10} = 6\)
  \item \(\frac{4z + 2}{6} + \frac{z - 1}{3} = 5\)
\end{enumerate}

\section*{Gemischte Aufgaben}

\subsection*{Aufgabe 9: Alles kombinieren}
Löse die Gleichungen. Nutze alle Methoden.
\begin{enumerate}
  \item \(\frac{x}{2} + \frac{x}{4} = 6\)
  \item \(\frac{3y}{5} + 2 = \frac{y}{3} + 7\)
  \item \(\frac{2a + 1}{4} - \frac{a}{8} = 3\)
  \item \(\frac{5z - 2}{10} + \frac{z}{5} = 8\)
\end{enumerate}

\subsection*{Aufgabe 10: Prüfen}
Setze deine Lösungen in die ursprünglichen Gleichungen ein und prüfe, ob sie stimmen.

\end{document}
